\documentclass[12pt,a4paper]{article}

% ── Packages ─────────────────────────────────────────────────────────────────
\usepackage[margin=2.5cm]{geometry}
\usepackage{booktabs}
\usepackage{graphicx}
\usepackage{amsmath}
\usepackage{hyperref}
\usepackage{parskip}
\usepackage{caption}
\usepackage{microtype}
\usepackage[numbers]{natbib}
\usepackage{setspace}
\onehalfspacing

\hypersetup{colorlinks=true, linkcolor=blue, citecolor=blue, urlcolor=blue}

% ── Title ─────────────────────────────────────────────────────────────────────
\title{%
  \textbf{Learning-Augmented Page Replacement Under Workload Shifts}\\[6pt]
  \large CSE-307: Operating Systems --- Term Paper (Track 1)
}
\author{%
  \textbf{Tanvir Ahmed Rafi}\\
  Student ID: 202414025 \quad\textbar\quad Section: A \quad\textbar\quad Batch: CSE-24\\
  Department of Computer Science and Engineering\\[4pt]
  \small\textbf{GitHub Repository:} \url{https://github.com/Tanvir-Ahmed-Rafi/cse307-learned-page-replacement}
}
\date{October 2026}

\begin{document}
\maketitle
\thispagestyle{empty}

% ─────────────────────────────────────────────────────────────────────────────
\section{Problem Framing}
% ─────────────────────────────────────────────────────────────────────────────

Modern operating systems use \emph{page replacement} to decide which physical
memory page to evict when a page fault occurs and all frames are full
\cite{silberschatz2018}. Classical algorithms such as FIFO and LRU perform
reasonably under stationary workloads but lack the ability to adapt when the
access pattern changes abruptly. Real programs, however, routinely shift
between phases: a startup phase that exhibits strong locality, followed by
a streaming or scanning phase that accesses pages almost uniformly at random.

This paper investigates whether a lightweight machine-learning policy---a
Decision Tree classifier---can narrow the gap with Belady's Optimal algorithm
\cite{belady1966} under such a two-phase workload shift. The central questions
are: (1) how much does each algorithm degrade after the shift, and (2) does
training on the locality phase help or hurt the learned policy during the
random phase?

% ─────────────────────────────────────────────────────────────────────────────
\section{Implementation Summary}
% ─────────────────────────────────────────────────────────────────────────────

All algorithms are implemented from scratch in Python 3.

\paragraph{FIFO.} A queue records the insertion order of pages. On eviction,
the page at the front of the queue (oldest arrival) is removed. FIFO is
simple but ignores how recently a page has been used \cite{tanenbaum2015}.

\paragraph{LRU.} Each page carries a timestamp of its most recent access.
On eviction, the page with the smallest timestamp (least recently used) is
removed. LRU exploits temporal locality: recently used pages are likely to be
needed again soon \cite{silberschatz2018}.

\paragraph{Optimal (Belady's).} On eviction, the page whose next use is
farthest in the future is removed. If a page is never accessed again, it is
evicted first. Because future access information is required, Optimal cannot
be deployed in a real OS, but it provides a provable lower bound on page faults
\cite{belady1966}. In this project Optimal also supplies ground-truth eviction
labels for training the learned policy.

\paragraph{Learned Policy.} A \texttt{DecisionTreeClassifier} (scikit-learn
\cite{pedregosa2011}, \texttt{max\_depth=4}, \texttt{min\_samples\_leaf=5}) is
trained to predict which page currently in memory is the best eviction
candidate. Three causal features are computed from past accesses only:

\begin{itemize}
  \item \textbf{recency} --- steps elapsed since the page was last accessed;
  \item \textbf{frequency} --- total number of accesses to the page so far;
  \item \textbf{recency\_ratio} --- \(\text{recency} / (\text{step} + 1)\),
        a normalised recency signal.
\end{itemize}

Future information is used exclusively to compute binary labels
(1 = Optimal would evict this page, 0 = otherwise) during training.
At inference time the model scores every candidate in memory and evicts the
one with the highest predicted probability of being the optimal target.

% ─────────────────────────────────────────────────────────────────────────────
\section{Experimental Setup}
% ─────────────────────────────────────────────────────────────────────────────

\subsection{Synthetic Trace}

A reproducible two-phase trace of 1{,}000 page accesses was generated using
\texttt{numpy.random.default\_rng(seed=42)}. Pages are integers in
$[0, 24]$ (25 distinct pages).

\begin{itemize}
  \item \textbf{Phase 1 (indices 0--499):} A hot working set of six pages
        (0--5) is accessed repeatedly. Approximately 40\% of accesses form
        short sequential bursts of length four (e.g.\ pages 2$\to$3$\to$4$\to$5),
        mimicking spatial locality. 10\% of accesses go to a background
        page outside the hot set.
  \item \textbf{Workload shift at index 500.}
  \item \textbf{Phase 2 (indices 500--999):} Page numbers are drawn
        uniformly at random from all 25 pages. No working-set discipline;
        every page is equally probable.
\end{itemize}

\subsection{Model Training}

Training samples were collected by replaying Phase~1 with 5~frames and
recording features at every eviction point, labelled using Belady's
next-use table. This yielded 425 samples (85 positive, 340 negative).
The data was split 75/25 into training (318 samples) and evaluation
(107 samples) sets using stratified sampling.

\subsection{Evaluation Protocol}

The same 1{,}000-access trace was used for all four algorithms.
Frame counts of 3, 4, 5, and 6 were tested. Statistics were collected
both for the full trace and separately for each phase.

% ─────────────────────────────────────────────────────────────────────────────
\section{Results}
% ─────────────────────────────────────────────────────────────────────────────

\subsection{Model Evaluation}

Table~\ref{tab:model} shows the Decision Tree's performance on the held-out
evaluation set. High precision (1.00) but low recall (0.33) indicates the
model is conservative: when it predicts eviction it is always correct, but
it misses two-thirds of the truly optimal evictions.

\begin{table}[h]
\centering
\caption{Decision Tree evaluation metrics (held-out test set, $n=107$)}
\label{tab:model}
\begin{tabular}{lc}
\toprule
Metric & Value \\
\midrule
Accuracy  & 0.8692 \\
Precision & 1.0000 \\
Recall    & 0.3333 \\
F1        & 0.5000 \\
\bottomrule
\end{tabular}
\end{table}

\subsection{Overall Page Fault Comparison}

Table~\ref{tab:overall} reports total page faults across all 1{,}000 accesses
for each algorithm and frame count. Optimal consistently achieves the fewest
faults. The Learned policy sits between Optimal and FIFO/LRU for all frame
counts, demonstrating that ML guidance provides a measurable benefit even with
simple features.

\begin{table}[h]
\centering
\caption{Total page faults (1{,}000 accesses)}
\label{tab:overall}
\begin{tabular}{lcccc}
\toprule
Algorithm & Frames = 3 & Frames = 4 & Frames = 5 & Frames = 6 \\
\midrule
FIFO    & 807 & 710 & 599 & 509 \\
LRU     & 810 & 718 & 612 & 476 \\
Optimal & 579 & 460 & 368 & 296 \\
Learned & 757 & 657 & 539 & 443 \\
\bottomrule
\end{tabular}
\end{table}

\subsection{Hit Ratios by Phase (Frames = 4)}

Table~\ref{tab:phase} shows hit ratios split across the two phases for
$\text{frames}=4$ (the representative case). Every algorithm degrades
substantially after the workload shift.

\begin{table}[h]
\centering
\caption{Hit ratio before and after workload shift (frames = 4)}
\label{tab:phase}
\begin{tabular}{lccc}
\toprule
Algorithm & Phase 1 & Phase 2 & Drop \\
\midrule
FIFO    & 0.432 & 0.148 & $-$0.284 \\
LRU     & 0.412 & 0.152 & $-$0.260 \\
Optimal & 0.698 & 0.382 & $-$0.316 \\
Learned & 0.546 & 0.140 & $-$0.406 \\
\bottomrule
\end{tabular}
\end{table}

% ─────────────────────────────────────────────────────────────────────────────
\section{Analysis of Workload Shift}
% ─────────────────────────────────────────────────────────────────────────────

All algorithms experience a dramatic drop in hit ratio after the shift.
Several observations are noteworthy.

\textbf{Why does performance drop?} In Phase~2 accesses are spread uniformly
across 25 pages while only 4 frames are available, giving a theoretical maximum
hit ratio of $4/25 = 0.16$. No algorithm can maintain the high hit ratios seen
in Phase~1 because the working-set size now exceeds the frame count.

\textbf{Why does Optimal retain the highest hit ratio in both phases?}
Optimal always evicts the page that will not be needed for the longest time.
Even in a random trace this is better than any heuristic, because any page
that happens to recur quickly is kept in memory.

\textbf{Why does the Learned policy degrade the most (drop of 0.406)?}
The Decision Tree was trained exclusively on Phase-1 data, so it learned that
pages with high recency (stale pages in the locality-heavy hot set) should be
evicted. In Phase~2 there is no stable hot set; the features the model relies
on no longer carry the same predictive signal. The model therefore makes
suboptimal evictions, performing slightly worse than LRU ($-0.406$ vs.\
$-0.260$) in Phase~2. This illustrates \emph{distribution shift} in an OS
context: a policy trained on one phase can hurt performance in a different phase.

\textbf{Interesting FIFO vs.\ LRU reversal.} At frames = 3, 4, 5 FIFO has
fewer faults than LRU; at frames = 6 LRU outperforms FIFO. This kind of
frame-count sensitivity is consistent with the Belady anomaly and the known
fragility of FIFO under varying frame counts \cite{tanenbaum2015}.

% ─────────────────────────────────────────────────────────────────────────────
\section{Conclusion}
% ─────────────────────────────────────────────────────────────────────────────

This project shows that a lightweight Decision Tree, trained on Belady's
optimal labels with three causal features, can outperform both FIFO and LRU
during the training phase (Phase~1 hit ratio: 0.546 vs.\ 0.432 for FIFO) but
suffers from distribution shift during the unseen random phase (Phase~2 hit
ratio: 0.140, slightly below LRU at 0.152). The gap between the Learned policy
and Optimal remains large (0.546 vs.\ 0.698 in Phase~1), primarily because the
model's low recall (0.33) means it misses two-thirds of the truly best evictions.

Future directions include online model updates (retrain on a sliding window),
richer features (page-fault rate, inter-access time distribution), and
evaluation on multi-phase real-world memory traces. Despite its limitations,
this experiment demonstrates that even a very simple ML model can add
measurable value over classical heuristics when trained on representative data.

% ─────────────────────────────────────────────────────────────────────────────
\section*{References}
% ─────────────────────────────────────────────────────────────────────────────
\begingroup
\renewcommand{\section}[2]{}%  suppress duplicate heading
\bibliographystyle{plainnat}

\begin{thebibliography}{9}

\bibitem{silberschatz2018}
Silberschatz, A., Galvin, P.~B., and Gagne, G.\ (2018).
\newblock \emph{Operating System Concepts}, 10th ed.
\newblock Wiley.

\bibitem{belady1966}
Belady, L.~A.\ (1966).
\newblock A study of replacement algorithms for a virtual-storage computer.
\newblock \emph{IBM Systems Journal}, 5(2):78--101.

\bibitem{tanenbaum2015}
Tanenbaum, A.~S., and Bos, H.\ (2015).
\newblock \emph{Modern Operating Systems}, 4th ed.
\newblock Pearson.

\bibitem{pedregosa2011}
Pedregosa, F., et al.\ (2011).
\newblock Scikit-learn: Machine learning in Python.
\newblock \emph{Journal of Machine Learning Research}, 12:2825--2830.

\end{thebibliography}
\endgroup

\end{document}
